% GENERATED FILE. Edit canonical lessons, then run npm run generate:books.
% canonical-source-hash: fnv1a64:2eff6297a9082cb4
% canonical-lessons: SA-C453-utkantha, SA-C453-vidhurah, SA-C453-pravasati, SA-C453-suhrd, SA-C453-sandisati

\chapter{Missing Someone}
\label{ch:sa-missing-someone}
\hlchaptermodality{453}

\begin{chapteropening}
\textbf{By the end of this chapter:} \emph{I can say that I miss someone far away.}

Complete the last lesson of chapter 453: I can say that I miss someone far away.
\end{chapteropening}

\section[\texorpdfstring{\sk{उत्कण्ठा} (utkaṇṭhā)}{उत्कण्ठा (utkaṇṭhā)}]{\sk{उत्कण्ठा} (utkaṇṭhā) — pining for someone absent}

\label{lesson:SA-C453-utkantha}



\begin{quote}
Before the new one: say the Sanskrit for agitated, then the Sanskrit for wretched.
\end{quote}

\subsection*{You'll want to know: \sk{उत्कण्ठा}}
\textbf{\sk{उत्कण्ठा}} — \emph{utkaṇṭhā} — \textquotedblleft{}pining for someone absent\textquotedblright{}.

\begin{quote}
\sk{प्रवासे} \sk{मम} \sk{गृहस्य} \sk{उत्कण्ठा} \sk{भवति।}
\end{quote}

\begin{quote}
\emph{pravāse mama gṛhasya utkaṇṭhā bhavati.}
\end{quote}

\begin{quote}
When I am away, I long for home.
\end{quote}

A feminine noun.

\begin{cousinweb}
\textbf{\sk{उत्कण्ठा}} (\emph{utkaṇṭhā}) is built from \emph{ud}, \textquotedblleft{}up\textquotedblright{}, and \emph{kaṇṭha}, \textquotedblleft{}the neck\textquotedblright{} (\textbf{\sk{कण्ठः}}): MW's \emph{ut-kaṇṭha} is \textquotedblleft{}having the neck uplifted\textquotedblright{}, and \emph{utkaṇṭhā} \textquotedblleft{}longing for a beloved person, missing someone\textquotedblright{} (MW).
\end{cousinweb}

\begin{grammarlens}[title={What you've built}]
One more word for describing what you lived and felt.
\end{grammarlens}

\subsection*{Your turn}
\begin{itemize}
\raggedright
  \item \emph{Say it:} \emph{utkaṇṭhā}
  \item \emph{Say it:} \emph{utkaṇṭhā}, once more
  \item \emph{Say it:} \emph{pravāse mama gṛhasya utkaṇṭhā bhavati.} — When I am away, I long for home.
  \item \emph{Recall:} say the Sanskrit for agitated, then the Sanskrit for wretched, then say \emph{utkaṇṭhā} again
\end{itemize}

\begin{tcolorbox}[breakable,colback=teal!4,colframe=teal!35!black,title={Before you move on}]
What does \sk{उत्कण्ठा} mean? (\textquotedblleft{}Pining for someone absent\textquotedblright{}.) Say it once more.
\end{tcolorbox}
\section[\texorpdfstring{\sk{विधुरः} (vidhuraḥ)}{विधुरः (vidhuraḥ)}]{\sk{विधुरः} (vidhuraḥ) — bereft; alone}

\label{lesson:SA-C453-vidhurah}



\begin{quote}
Before the new one: say the Sanskrit for wretched, then the Sanskrit for pining for someone absent.
\end{quote}

\subsection*{You'll want to know: \sk{विधुरः}}
\textbf{\sk{विधुरः}} — \emph{vidhuraḥ} — \textquotedblleft{}bereft; alone\textquotedblright{}.

\begin{quote}
\sk{मित्रं} \sk{दूरे} \sk{वसति}, \sk{अतः} \sk{अहं} \sk{विधुरः} \sk{अस्मि।}
\end{quote}

\begin{quote}
\emph{mitraṃ dūre vasati, ataḥ ahaṃ vidhuraḥ asmi.}
\end{quote}

\begin{quote}
My friend lives far away, so I feel alone.
\end{quote}

\begin{cousinweb}
\textbf{\sk{विधुरः}} (\emph{vidhuraḥ}) MW gives \textquotedblleft{}bereft, bereaved (esp. of any loved person), alone, solitary\textquotedblright{}; MW gives no derivation, so its origin is uncertain (MW).
\end{cousinweb}

\begin{grammarlens}[title={What you've built}]
One more word for describing what you lived and felt.
\end{grammarlens}

\subsection*{Your turn}
\begin{itemize}
\raggedright
  \item \emph{Say it:} \emph{vidhuraḥ}
  \item \emph{Say it:} \emph{vidhuraḥ}, once more
  \item \emph{Say it:} \emph{mitraṃ dūre vasati, ataḥ ahaṃ vidhuraḥ asmi.} — My friend lives far away, so I feel alone.
  \item \emph{Recall:} say the Sanskrit for wretched, then the Sanskrit for pining for someone absent, then say \emph{vidhuraḥ} again
\end{itemize}

\begin{tcolorbox}[breakable,colback=teal!4,colframe=teal!35!black,title={Before you move on}]
What does \sk{विधुरः} mean? (\textquotedblleft{}Bereft; alone\textquotedblright{}.) Say it once more.
\end{tcolorbox}
\section[\texorpdfstring{\sk{प्रवसति} (pravasati)}{प्रवसति (pravasati)}]{\sk{प्रवसति} (pravasati) — lives away from home}

\label{lesson:SA-C453-pravasati}



\begin{quote}
Before the new one: say the Sanskrit for pining for someone absent, then the Sanskrit for bereft.
\end{quote}

\subsection*{You'll want to know: \sk{प्रवसति}}
\textbf{\sk{प्रवसति}} — \emph{pravasati} — \textquotedblleft{}lives away from home\textquotedblright{}.

\begin{quote}
\sk{मम} \sk{भ्राता} \sk{दूरे} \sk{नगरे} \sk{प्रवसति।}
\end{quote}

\begin{quote}
\emph{mama bhrātā dūre nagare pravasati.}
\end{quote}

\begin{quote}
My brother lives away from home in a distant city.
\end{quote}

Past (he …): \sk{सः} \sk{संवत्सरं} \sk{प्रोषितवान्}, \emph{saḥ saṃvatsaraṃ proṣitavān}.

\begin{cousinweb}
\textbf{\sk{प्रवसति}} (\emph{pravasati}) is built from \emph{pra}, \textquotedblleft{}forth\textquotedblright{}, and the root √vas, \textquotedblleft{}to dwell\textquotedblright{}, the root of \textbf{\sk{वसति}}; MW gives \textquotedblleft{}to go or sojourn abroad\textquotedblright{} (MW).
\end{cousinweb}

\begin{grammarlens}[title={What you've built}]
One more word for describing what you lived and felt.
\end{grammarlens}

\subsection*{Your turn}
\begin{itemize}
\raggedright
  \item \emph{Say it:} \emph{pravasati}
  \item \emph{Say it:} \emph{pravasati}, once more
  \item \emph{Say it:} \emph{mama bhrātā dūre nagare pravasati.} — My brother lives away from home in a distant city.
  \item \emph{Recall:} say the Sanskrit for pining for someone absent, then the Sanskrit for bereft, then say \emph{pravasati} again
\end{itemize}

\begin{tcolorbox}[breakable,colback=teal!4,colframe=teal!35!black,title={Before you move on}]
What does \sk{प्रवसति} mean? (\textquotedblleft{}Lives away from home\textquotedblright{}.) Say it once more.
\end{tcolorbox}
\section[\texorpdfstring{\sk{सुहृद्} (suhṛd)}{सुहृद् (suhṛd)}]{\sk{सुहृद्} (suhṛd) — a dear friend; a well-wisher}

\label{lesson:SA-C453-suhrd}



\begin{quote}
Before the new one: say the Sanskrit for bereft, then the Sanskrit for lives away from home.
\end{quote}

\subsection*{You'll want to know: \sk{सुहृद्}}
\textbf{\sk{सुहृद्}} — \emph{suhṛd} — \textquotedblleft{}a dear friend; a well-wisher\textquotedblright{}.

\begin{quote}
\sk{सुहृद्} \sk{दुःखे} \sk{अपि} \sk{मित्रं} \sk{न} \sk{त्यजति।}
\end{quote}

\begin{quote}
\emph{suhṛd duḥkhe api mitraṃ na tyajati.}
\end{quote}

\begin{quote}
A true friend does not abandon you even in trouble.
\end{quote}

A masculine noun.

\begin{cousinweb}
\textbf{\sk{सुहृद्}} (\emph{suhṛd}) is built from \emph{su}, \textquotedblleft{}good\textquotedblright{}, and \emph{hṛd}, \textquotedblleft{}heart\textquotedblright{}: MW gives \textquotedblleft{}good-hearted, well-disposed; a friend, ally\textquotedblright{} (MW).
\end{cousinweb}

\begin{grammarlens}[title={What you've built}]
One more word for describing what you lived and felt.
\end{grammarlens}

\subsection*{Your turn}
\begin{itemize}
\raggedright
  \item \emph{Say it:} \emph{suhṛd}
  \item \emph{Say it:} \emph{suhṛd}, once more
  \item \emph{Say it:} \emph{suhṛd duḥkhe api mitraṃ na tyajati.} — A true friend does not abandon you even in trouble.
  \item \emph{Recall:} say the Sanskrit for bereft, then the Sanskrit for lives away from home, then say \emph{suhṛd} again
\end{itemize}

\begin{tcolorbox}[breakable,colback=teal!4,colframe=teal!35!black,title={Before you move on}]
What does \sk{सुहृद्} mean? (\textquotedblleft{}A dear friend; a well-wisher\textquotedblright{}.) Say it once more.
\end{tcolorbox}
\section[\texorpdfstring{\sk{सन्दिशति} (sandiśati)}{सन्दिशति (sandiśati)}]{\sk{सन्दिशति} (sandiśati) — sends word; sends a message}

\label{lesson:SA-C453-sandisati}



\begin{quote}
Before the new one: say the Sanskrit for lives away from home, then the Sanskrit for a dear friend.
\end{quote}

\subsection*{You'll want to know: \sk{सन्दिशति}}
\textbf{\sk{सन्दिशति}} — \emph{sandiśati} — \textquotedblleft{}sends word; sends a message\textquotedblright{}.

\begin{quote}
\sk{भ्राता} \sk{नगरात्} \sk{प्रतिदिनं} \sk{सन्दिशति।}
\end{quote}

\begin{quote}
\emph{bhrātā nagarāt pratidinaṃ sandiśati.}
\end{quote}

\begin{quote}
My brother sends word from the city every day.
\end{quote}

Past (he …): \sk{सः} \sk{नगरात्} \sk{सन्दिष्टवान्}, \emph{saḥ nagarāt sandiṣṭavān}.

\begin{cousinweb}
\textbf{\sk{सन्दिशति}} (\emph{sandiśati}) is built from \emph{sam}, \textquotedblleft{}together\textquotedblright{}, and the root √diś, \textquotedblleft{}to point out, to show\textquotedblright{}; MW gives \textquotedblleft{}to state, tell, direct\textquotedblright{} and \textquotedblleft{}to give an order or message to\textquotedblright{} someone (MW).
\end{cousinweb}

\begin{grammarlens}[title={What you've built}]
One more word for describing what you lived and felt.
\end{grammarlens}

\subsection*{Your turn}
\begin{itemize}
\raggedright
  \item \emph{Say it:} \emph{sandiśati}
  \item \emph{Say it:} \emph{sandiśati}, once more
  \item \emph{Say it:} \emph{bhrātā nagarāt pratidinaṃ sandiśati.} — My brother sends word from the city every day.
  \item \emph{Recall:} say the Sanskrit for lives away from home, then the Sanskrit for a dear friend, then say \emph{sandiśati} again
\end{itemize}

\begin{tcolorbox}[breakable,colback=teal!4,colframe=teal!35!black,title={Before you move on}]
What does \sk{सन्दिशति} mean? (\textquotedblleft{}Sends word; sends a message\textquotedblright{}.) Say it once more.
\end{tcolorbox}
